\chapter{Frequency-domain ROMs, greedy training and certified bounds}\label{ch:freq}

\begin{tcolorbox}[colback=hpale,colframe=hnavy!60,coltitle=white,colbacktitle=hnavy!85,title={In this chapter}]
Sweep a damped beam over frequency and the response is flat for long stretches and then jumps at each resonance. A reduced model needs full-order snapshots, each one an expensive solve, and the question is where to take them: evenly spaced, or where the response is hardest. We build the frequency-domain machinery step by step. The matrix $\bm A(\omega)=-\omega^2\bm M+i\omega\bm C+\bm K$ is nearly singular at every resonance, so a greedy loop is used to choose frequencies; an error indicator or a certified bound is used to decide how far to trust the result without solving the full model; and the reduced model is used for a random load described by a power spectral density. Everything is measured against the full-order frequency solver of \fea{}, and the certified bound, the one rigorous error statement in the package, is tested for how conservative it is.\par\smallskip
After it you should be able to train a frequency-domain model by greedy selection, tell a cheap indicator from a certified bound, and compute a random-load response from a reduced model.
\end{tcolorbox}

\section{The problem}\label{sec:fr-problem}

Suppose we have a budget of five full-order solves to build a reduced model of a cantilever over a band that contains three resonances. The obvious plan is to spread the five evenly over the band. It is a fair baseline, but the beam does not care about even spacing: its response changes fast near resonances and slowly between them, so evenly spaced snapshots may miss exactly the places that matter. This chapter asks whether the loop can choose better, and how we would know. The beam of the example script is the running example.

\begin{sidenote}{Frequency-domain response and complex unknowns}
In the frequency domain a harmonic load $\bm F e^{i\omega t}$ produces a response $\bm x(\omega)e^{i\omega t}$ with $\bm x(\omega)$ complex, since damping shifts its phase. The unknown is therefore a complex vector at each frequency, obtained from one linear solve of the dynamic stiffness $\bm A(\omega)=-\omega^2\bm M+i\omega\bm C+\bm K$. A sweep over many frequencies repeats the solve, which is the cost the reduced model is meant to avoid.
\end{sidenote}

For harmonic excitation $\bm f(t)=\bm F\,e^{i\omega t}$ the steady response of $\bm M\ddot{\bm x}+\bm C\dot{\bm x}+\bm K\bm x=\bm f$ is $\bm x(\omega)e^{i\omega t}$ with
\begin{equation}\label{eq:fr-fom}
\bm A(\omega)\,\bm x(\omega)=\bm F,\qquad \bm A(\omega)=-\omega^2\bm M+i\omega\bm C+\bm K .
\end{equation}
A frequency sweep solves \eqref{eq:fr-fom} for hundreds or thousands of frequencies. Near each natural frequency $\bm A(\omega)$ is almost singular and $\bm x(\omega)$ changes sharply, so a basis has to resolve the resonances in the band, and any measure of error that divides by a small quantity is delicate there.

\paragraph{The model problem.}
A damped cantilever of $20$ Euler--Bernoulli elements (\cfrNFree{} free unknowns), Rayleigh damping $\bm C=\alpha\bm M+\beta\bm K$ with $\alpha=2$, $\beta=10^{-5}$, a unit tip-force pattern scaled to $1000$~N, the model of \path{rom_engine/examples/greedy_frequency_training.py}. The first natural frequency is \cfrFOne{}~Hz. The band is $0.1\omega_1$ to $20\omega_1$ and contains the first three resonances, the third at $\cfrThirdRes\,\omega_1$. The reference is \fea{}'s own harmonic solver, \code{FESystem.solve_harmonic}.

\section{Mathematical structure}\label{sec:fr-math}

\paragraph{Affine in $\omega$.}
The operator $\bm A(\omega)$ is a sum of fixed matrices with scalar coefficients, so the machinery of Chapter~\ref{ch:pod} applies with $\omega$ as the parameter. For Rayleigh damping the three terms collapse to two,
\begin{equation}\label{eq:fr-affine}
\bm A(\omega)=\theta_1(\omega)\bm M+\theta_2(\omega)\bm K,\qquad \theta_1=-\omega^2+i\omega\alpha,\quad \theta_2=1+i\omega\beta,
\end{equation}
and for general damping $\bm A=-\omega^2\bm M+i\omega\bm C+\bm K$ has three. \anchor{rom_engine.frequency.FrequencyROM} builds the matching \anchor{rom_engine.affine.AffineDecomposition} and projects it on a basis $\bm V$; each query is a small complex solve $\bm A_r(\omega)\bm q=\bm V^H\bm F$.

\paragraph{Which basis.}
A modal basis (the lowest undamped modes) does not depend on the load; it represents the response only as well as the load excites the retained modes and the discarded modes' static flexibility is small. A basis from the POD of FRF snapshots (full-order solves at selected frequencies, complex-valued) is built from the actual response to this load. Both are Galerkin projections, so the quality is the quality of the subspace.

\begin{sidenote}{Snapshot and column}
A snapshot is one full-order solution, here the complex response at one frequency. The basis is the POD of the snapshots, and each of its columns is a complex vector. Counting columns is the fair way to compare bases of equal size, but the modal columns are real, so a complex column carries twice the numbers.
\end{sidenote}

\begin{sidenote}{The greedy idea}
A \emph{greedy} algorithm makes the best choice available at each step without looking ahead. Here the choice is where to put the next full-order solve: at the frequency at which the current reduced model is judged worst. A new snapshot there is added to the basis, the reduced model is rebuilt, and the process repeats until the worst-case estimate is small enough. The result is a basis tuned to the frequencies where it is needed, at a cost of one full solve per added column.
\end{sidenote}

\paragraph{Greedy training.}
The expensive part is the full-order solves that produce the snapshots. A weak-greedy loop starts from a few seed frequencies, builds a reduced model, evaluates a cheap error \emph{indicator} at every candidate frequency in a fine pool, solves the full model at the candidate where the indicator is largest, and repeats. Two indicators are available: the residual norm $\|\bm F-\bm A(\omega)\bm V\bm q\|$, evaluated in $O(Qnr)$ per frequency, and the hierarchical indicator $\|\bm x_r(\omega)-\bm x_{r'}(\omega)\|$, the disagreement of the current reduced model with the previous one. Neither is a bound. The first is a residual, the second assumes that the error keeps shrinking as snapshots are added.

\begin{sidenote}{Indicator and bound}
An indicator ranks frequencies by how bad the reduced model probably is, and is used to choose where to look next. A certified bound is a number that the true error cannot exceed. The first is cheap and says nothing about size; the second is rigorous and may be so loose that it says little.
\end{sidenote}

\paragraph{A certified bound.}
For any $\bm x_r$, $\bm A(\bm x-\bm x_r)=\bm F-\bm A\bm x_r=:\bm\rho$, so
\begin{equation}\label{eq:fr-bound}
\|\bm x(\omega)-\bm x_r(\omega)\|_2\;\le\;\frac{\|\bm\rho(\omega)\|_2}{\sigma_{\min}(\bm A(\omega))}\;\le\;\frac{\|\bm\rho(\omega)\|_2}{\beta_{\mathrm{LB}}(\omega)},
\end{equation}
where $\beta_{\mathrm{LB}}\le\sigma_{\min}$ is any lower bound that can be evaluated without the full model; this is the standard a posteriori estimate of the reduced-basis method~\cite{hesthaven2016rb}. Computing $\sigma_{\min}$ at every $\omega$ would need a full-size singular value decomposition, so the lower bound is built offline from a few exact values at reference frequencies. The package has two constructions. The simple one uses Weyl's inequality $|\sigma_{\min}(\bm A)-\sigma_{\min}(\bm B)|\le\|\bm A-\bm B\|_2$ and the affine form: $\sigma_{\min}(\bm A(\omega))\ge\sigma_{\min}(\bm A(\omega_j))-\sum_q|\theta_q(\omega)-\theta_q(\omega_j)|\,\|\bm A_q\|_2$, maximised over references $\omega_j$. The second is the successive constraint method of Huynh, Rozza, Sen and Patera~\cite{huynh2007scm}, applied to $\bm A^H\bm A$, which is real, symmetric and positive semidefinite for any complex $\bm A$; it answers each query by a small linear program over the reference data. Both are rigorous whenever the lower bound is positive, and both return $+\infty$ when it is not.

\paragraph{Random vibration.}
For a stationary random load with a fixed spatial pattern $\bm F_0$ and input spectrum $S_{\mathrm{in}}(f)$, the output PSD at a sensor is $S_{\mathrm{out}}(f)=|H(f)|^2S_{\mathrm{in}}(f)$ with $H$ the transfer function from $\bm F_0$ to that output, and the RMS response is $\sigma=\bigl(\int S_{\mathrm{out}}\,df\bigr)^{1/2}$. A reduced model replaces $H$ by $H_r$; the integral is evaluated by the trapezoidal rule on the frequency grid.

\begin{algbox}{PSD response and RMS from a reduced model}
\begin{pseudo}
\Require reduced model $H_r$, input spectrum $S_{\mathrm{in}}(f)$, frequency grid $f_1<\dots<f_m$
\For{$i=1,\dots,m$}
  \State $H_r(f_i)\gets$ small complex solve for the sensor output;\ $S_{\mathrm{out}}(f_i)\gets|H_r(f_i)|^2S_{\mathrm{in}}(f_i)$
\EndFor
\State $\sigma\gets\big(\sum_i\tfrac12(S_{\mathrm{out},i}+S_{\mathrm{out},i+1})\,\Delta f_i\big)^{1/2}$ \Comment{trapezoidal rule}
\Ensure $S_{\mathrm{out}}(f)$ and RMS $\sigma$
\end{pseudo}
\end{algbox}

The cost is $m$ small solves, where the full-order model would need $m$ full ones. The grid has to resolve the peaks, with several points per half-power width for lightly damped modes, and the load has to be a single fixed spatial pattern. The code is \anchor{rom_engine.random_vibration.psd_response}; \Cref{tab:fr-psd} is the check.

\begin{algbox}{Certified error bound at a frequency}
\begin{pseudo}
\Require reduced solution $\bm x_r(\omega)$, references $\omega_1,\dots,\omega_J$ with exact $\sigma_{\min}(\bm A(\omega_j))$ (offline), affine terms $\theta_q$, $\|\bm A_q\|_2$
\State $\bm\rho\gets\bm F-\bm A(\omega)\bm x_r$ \Comment{full-size residual, one product per matrix term}
\State $\beta_{\mathrm{LB}}\gets\max_j\big[\sigma_{\min}(\bm A(\omega_j))-\sum_q|\theta_q(\omega)-\theta_q(\omega_j)|\,\|\bm A_q\|_2\big]$ \Comment{Weyl; or an SCM linear program}
\If{$\beta_{\mathrm{LB}}>0$} $\Delta\gets\|\bm\rho\|/\beta_{\mathrm{LB}}$ \Else{} $\Delta\gets+\infty$ \EndIf
\Ensure $\|\bm x(\omega)-\bm x_r(\omega)\|\le\Delta$
\end{pseudo}
\end{algbox}

Online the cost is one full-size residual and $O(J\,Q)$ work; offline it is $J$ exact smallest singular values of full-size matrices. Near a resonance the lower bound can fall to zero or below, and the certified bound is then $+\infty$, which is correct and useless; the references can also be too sparse for the affine terms to move far from them. The code is \anchor{rom_engine.scm.SingularValueLowerBound}, \anchor{rom_engine.scm_lp.LPSingularValueLowerBound} and \anchor{rom_engine.scm.certified_error_bound}; \Cref{fig:fr-bound} compares the bound with the true error.

\tikzfigure[0.9\linewidth]{ch10_bound}{The certified bound. The residual is computed online at full size; the lower bound on $\sigma_{\min}$ is built offline from a few exact values.}{fig:fr-boundidea}

\begin{algbox}{Greedy frequency-domain ROM with an optional certified bound}
\begin{pseudo}
\Require band, candidate pool $\Omega$, load pattern $\bm F$, seed frequencies, column budget or tolerance
\State full-order solves at the seeds;\ POD basis $\bm V$ of the complex snapshots;\ build $\bm A_r(\omega)$
\Repeat
  \State indicator $I(\omega)$ for every unused $\omega\in\Omega$ \Comment{residual norm in round 1, then hierarchical}
  \State $\omega^\ast\gets\arg\max I$;\ full-order solve at $\omega^\ast$;\ add the snapshot;\ rebuild $\bm V$
\Until{column budget reached or $\max I<\mathrm{tol}$}
\State optional: build $\beta_{\mathrm{LB}}$ and report $\|\bm\rho\|/\beta_{\mathrm{LB}}$ where finite; otherwise validate on frequencies near the resonances that were not used in training
\Ensure basis $\bm V$, list of selected $\omega^\ast$ (for audit)
\end{pseudo}
\end{algbox}

Offline the cost is one full-order solve per selected frequency plus $O(Qnr)$ per candidate and round; online it is one small complex solve per query. The loop goes wrong when the indicator is small away from a resonance that was never sampled, because neither indicator is a bound, and when the candidate pool does not contain the peak. The code is \anchor{rom_engine.greedy.greedy_train_frequency_basis} and \anchor{rom_engine.frequency.build_pod_basis_from_frf_snapshots}; \Cref{fig:fr-greedy}, \Cref{fig:fr-ind} and \Cref{tab:fr-greedy} are the checks.

\tikzfigure[0.98\linewidth]{ch10_greedy}{The greedy loop in the frequency domain. Full-order solves happen only at the frequencies chosen by the indicator.}{fig:fr-greedyidea}

\section{Implementation in the packages}\label{sec:fr-impl}

\begin{lstlisting}[caption={Greedy training, the reduced model, indicators, a certified bound and a random-vibration RMS (K, M, C are free-dof matrices; F the load pattern).},label={lst:fr-run}]
import numpy as np
from rom_engine import FrequencyROM, greedy_train_frequency_basis
from rom_engine.scm_lp import LPSingularValueLowerBound, certified_error_bound
from rom_engine.random_vibration import psd_response

band = np.linspace(0.1, 20.0, 400) * w1                 # candidate pool, rad/s
basis, history = greedy_train_frequency_basis(
    band, M, K, F, rayleigh=(alpha, beta), n_seed=3, tol=1e-12, max_modes=5)
print([om / w1 for om, _, _ in history])                # frequencies the loop chose
rom = FrequencyROM.from_MCK(M, K, basis.V, rayleigh=(alpha, beta))

print(rom.residual_norm(om, F))                         # indicator, not a bound
bound = LPSingularValueLowerBound.from_affine(rom.affine)
bound.greedy_train(candidates, tol=0.05, max_references=80)
print(certified_error_bound(rom, bound, om, F))         # inf where the lower bound is not positive

freqs = np.linspace(8.4, 1676.0, 6000)                  # Hz
S_out, rms = psd_response(rom, freqs, np.ones_like(freqs), F_full[free], tip)
\end{lstlisting}

\paragraph{Try it.} \Cref{lst:fr-try} spends the same five-column budget twice on the beam of the example script: evenly spaced snapshots first, then greedy selection. Both models are tested on $60$ frequencies scattered by $5\%$ around the first three resonances.

\begin{lstlisting}[caption={Five evenly spaced snapshots against five greedily chosen ones (run from \texttt{rom\_engine/examples}).},label={lst:fr-try}]
import numpy as np
from scipy.linalg import eigh
from greedy_frequency_training import build_damped_cantilever
from rom_engine import FrequencyROM, build_pod_basis_from_frf_snapshots, greedy_train_frequency_basis

sysobj, alpha, beta = build_damped_cantilever(n=20)
free, n_dof = sysobj.free_dofs, sysobj.n_dof
K, M, C = (a[np.ix_(free, free)] for a in (sysobj.K, sysobj.M, sysobj.C))
w1 = np.sqrt(eigh(K, M, eigvals_only=True)[0])
F = np.zeros(n_dof); F[n_dof - 2] = 1000.0
band = np.linspace(0.1, 20.0, 400) * w1

# naive: five snapshots at evenly spaced frequencies
uni = build_pod_basis_from_frf_snapshots(np.linspace(band[0], band[-1], 5), M, K, F[free],
                                         C=C, n_modes=5)
# greedy: five columns, the loop chooses where to solve
gre, history = greedy_train_frequency_basis(band, M, K, F[free], rayleigh=(alpha, beta),
                                            n_seed=3, tol=1e-12, max_modes=5)
print("greedy chose", [round(float(om / w1), 2) for om, _, _ in history], "x omega_1")

rng = np.random.default_rng(0)                   # 5% scatter around three resonances
test = np.concatenate([w * (1 + 0.05 * rng.standard_normal(20))
                       for w in np.sqrt(eigh(K, M, eigvals_only=True)[:3])])
for name, b in (("uniform", uni), ("greedy", gre)):
    rom = FrequencyROM.from_MCK(M, K, b.V, rayleigh=(alpha, beta))
    err = [np.linalg.norm(rom.frequency_response([w], F[free])[0] - sysobj.solve_harmonic(w, F)[free])
           / np.linalg.norm(sysobj.solve_harmonic(w, F)[free]) for w in test]
    print(f"{name:8s} 5 columns: median error {np.median(err):.1e}, max {np.max(err):.1e}")
\end{lstlisting}

\begin{lstlisting}[numbers=none,language={},keywordstyle={},commentstyle={},frame=single,aboveskip=2pt]
greedy chose [1.0, 7.18] x omega_1
uniform  5 columns: median error 1.6e-05, max 1.3e-04
greedy   5 columns: median error 3.0e-06, max 4.2e-04
\end{lstlisting}

The loop picks the first resonance and a frequency near the second without being told where they are, and its median error is five times lower than the uniform basis. Its maximum error is three times \emph{higher}. So greedy is not simply better, and the reason is the subject of the next section: the worst frequency sits at a resonance neither basis has seen.

\begin{table}[ht]
\centering\small
\begin{tabularx}{\linewidth}{@{}l X >{\centering\arraybackslash}p{0.11\linewidth}@{}}
\toprule
Concept & Code & Script calls \\
\midrule
reduced frequency-domain model, residual norm, hierarchical indicator & \anchor{rom_engine.frequency.FrequencyROM} & yes \\
POD of complex FRF snapshots & \anchor{rom_engine.frequency.build_pod_basis_from_frf_snapshots} & yes \\
weak-greedy selection of frequencies & \anchor{rom_engine.greedy.greedy_train_frequency_basis} & yes \\
Lipschitz lower bound on $\sigma_{\min}$ and certified error bound & \anchor{rom_engine.scm.SingularValueLowerBound}, \anchor{rom_engine.scm.certified_error_bound} & yes \\
successive constraint method (LP), with its own \code{certified_error_bound} & \anchor{rom_engine.scm_lp.LPSingularValueLowerBound} & yes \\
PSD response and RMS & \anchor{rom_engine.random_vibration.psd_response} & yes \\
band-limited Gaussian time history (not tested here) & \anchor{rom_engine.random_vibration.band_limited_gaussian_time_history} & no \\
worked examples & \path{rom_engine/examples/gallery/rom_03_greedy_frequency_training.py} & -- \\
\bottomrule
\end{tabularx}
\caption{Where the concepts of this chapter live in the repository (tag \booktag). The last column says whether this chapter's own checks call, by name, the code listed in the row: yes (every name), part, or no. Code reached only through another routine counts as no, so the column is a lower bound on what this book exercises.}\label{tab:fr-map}
\end{table}

\section{Verification}\label{sec:fr-verif}

All numbers were measured on \machine{}. Errors are relative, $\|\bm x_r-\bm x\|/\|\bm x\|$ against \code{solve_harmonic}. Two test sets, none of them used for training, are used: \cfrNTestNear{} frequencies scattered by $5\%$ around each of the first three resonances (the set of the package example), and \cfrNTestBand{} frequencies drawn uniformly over the band.

\subsection{Greedy against uniform against modal training}

\Cref{fig:fr-greedy} and \Cref{tab:fr-greedy} compare three bases of the same number of columns. A uniform basis uses full-order snapshots at evenly spaced frequencies, the greedy basis the loop above, and the modal basis the lowest undamped modes. The greedy loop starts from three evenly spaced seeds ($0.1$, $10.0$ and $20\,\omega_1$) and, for a basis of $5$ columns, picks two frequencies, \cfrSelFive{} in units of $\omega_1$: it lands on the first resonance and near the second ($6.27\,\omega_1$), without being told where they are, which is the behaviour the package documentation describes.

\begin{figure}[ht]
\centering
\includegraphics[width=0.98\linewidth]{figures/ch10_greedy.pdf}
\caption{Left: the full-order tip response and the two frequencies chosen by the greedy loop for a $5$-column basis (shaded). Right: maximum error on the test set near the resonances against the number of basis columns for the three bases.}\label{fig:fr-greedy}
\end{figure}

\begin{table}[ht]
\centering\small
\begin{tabular}{lrrrrrr}
\toprule
 & \multicolumn{3}{c}{5 columns} & \multicolumn{3}{c}{8 columns} \\
\cmidrule(lr){2-4}\cmidrule(lr){5-7}
basis & median & max (near) & max (band) & median & max (near) & max (band) \\
\midrule
greedy & \cfrGreedyMedFive & \cfrGreedyNearFive & \cfrGreedyBandFive & \cfrGreedyMedEight & \cfrGreedyNearEight & \cfrGreedyBandEight \\
uniform & \cfrUniformMedFive & \cfrUniformNearFive & \cfrUniformBandFive & \cfrUniformMedEight & \cfrUniformNearEight & \cfrUniformBandEight \\
modal & \cfrModalMedFive & \cfrModalNearFive & \cfrModalBandFive & \cfrModalMedEight & \cfrModalNearEight & \cfrModalBandEight \\
\bottomrule
\end{tabular}
\caption{Relative error of the three bases on frequencies not used for training. ``median'' and ``max (near)'' are over the \cfrNTestNear{} frequencies near the resonances, ``max (band)'' over the \cfrNTestBand{} band-wide frequencies.}\label{tab:fr-greedy}
\end{table}

The result does not support a general claim that greedy is better. At $5$ columns the greedy basis has a median error five times lower than the uniform one (\cfrGreedyMedFive{} against \cfrUniformMedFive{}) but a maximum error three times higher (\cfrGreedyNearFive{} against \cfrUniformNearFive{}); over the seven sizes tested ($3$ to $10$ columns) greedy has the lower median at \cfrGreedyWinsMed{} of \cfrRankCount{} and the lower maximum at \cfrGreedyWinsNear{} of \cfrRankCount{}. At $5$ columns the maximum of both bases sits at the third resonance (checked: the largest error is at $17.5\,\omega_1$ for each), which neither training covers, and it is the same peak in both test sets, which is why the near-resonance and band-wide maxima nearly coincide. Greedy selection puts the snapshots where the indicator is largest at the time, which is good for the resonances it finds and says nothing about the one it has not yet found. Both snapshot bases are far better than the modal one: at $5$ columns the modal basis has a maximum error of \cfrModalNearFive{}, since it ignores the load pattern, and greedy has the lower maximum at \cfrGreedyWinsModal{} of \cfrRankCount{} sizes. One caution: the snapshot bases are complex, so a ``column'' is a complex vector; the modal columns are real. This comparison is at equal column count, not at equal number of real degrees of freedom or equal offline cost (modal needs one eigensolve and no frequency solves).

\paragraph{Cost and local accuracy of the greedy $5$-column model.} The offline cost is the greedy training (full-order solves at the selected frequencies, basis construction) at \cfrTimeGreedyMs~ms plus the projection of the operator terms, below $0.1$~ms. At \cfrNFree{} unknowns one reduced solve takes \cfrTimeOnUs~$\mu$s and the lift to the full state \cfrTimeRecUs~$\mu$s, against \cfrTimeFullUs~$\mu$s for \fea{}'s harmonic solve and \cfrTimeDenseUs~$\mu$s for a plain dense solve of $\bm A(\omega)$. The speed-up of the reduced solve is a factor \cfrSpeedOnline{} against \fea{}'s solver, \cfrSpeedLift{} when the lift is included, and \cfrSpeedDense{} against the dense solve. These ratios exclude the offline cost. Including it, a sweep of \cfrWorkflowQ{} frequencies is \emph{not} faster than \cfrWorkflowQ{} full solves (a factor \cfrWorkflowSpeed), and the break-even is about \cfrBreakEven{} frequencies: at this size a frequency sweep is cheap enough that the full-order model is the better choice unless many more frequencies, or many load cases with the same basis, are needed. For a local quantity, the peak tip amplitude near each of the first three resonances (on a grid of $0.25\,\%$ steps in frequency) is reproduced to \cfrPeakAmpOne{}, \cfrPeakAmpTwo{} and \cfrPeakAmpThree{} relative error, and the frequency of the peak falls on the same grid point as the full model's in all three; this is much smaller than the largest error of the state over the near-resonance test set, \cfrPeakGlobal{}, which is taken over frequencies that are not at the peaks. Here the global state error is the more pessimistic of the two.

\subsection{Error indicators}

For the greedy $5$-column model, with the previous ($4$-column) model as the comparison, \Cref{fig:fr-ind} evaluates both indicators at \cfrNCand{} candidate frequencies against the true error. The residual norm ranks the frequencies in the same order as the true error (Spearman rank correlation \cfrSpearResid{}, and the \cfrTopK{} worst frequencies by residual are the \cfrTopK{} worst by true error, overlap \cfrTopOverlapResid{} of \cfrTopK{}). The hierarchical indicator ranks them less well (correlation \cfrSpearHier{}, overlap \cfrTopOverlapHier{} of \cfrTopK{}) and is not a measure of the current model's error at all: it measures the disagreement with the \emph{previous} model, which is larger than the current error by a factor between \cfrHierUnderMax{} and \cfrHierOverMax{} here. Neither gives the size of the error; each is a ranking signal for placing the next snapshot.

\begin{figure}[ht]
\centering
\includegraphics[width=0.98\linewidth]{figures/ch10_indicators.pdf}
\caption{The $5$-column greedy model on \cfrNCand{} candidate frequencies. Left: true error, residual norm and hierarchical indicator against frequency. Right: the two indicators against the true error.}\label{fig:fr-ind}
\end{figure}

\subsection{The certified bound}

With the greedy $5$-column model as the reduced model, both lower-bound constructions were trained with their own greedy reference selection (tolerance $0.05$, at most $15$ references) and tested on \cfrNGrid{} frequencies. The lower bounds were valid at every one of them (never above the true $\sigma_{\min}$, by an exact SVD). The certified error bound \eqref{eq:fr-bound} was finite at only \cfrScmBoundFinite{} of the \cfrNGrid{} frequencies for the Lipschitz construction and at \cfrLpBoundFinite{} for the LP; every finite value was above the true error (all of them, $1$ of $1$ and $8$ of $8$), and the bound overestimated the error by a factor of \cfrScmRatioMin{} (Lipschitz) and between \cfrLpRatioMin{} and \cfrLpRatioMax{} (LP). All \cfrScmInfNear{} grid frequencies within $5\%$ of a resonance have an infinite bound.

\Cref{tab:fr-sweep} varies the number of references. The Lipschitz bound stays finite at one frequency even with \cfrSweepLipMax{} references. The reason is in the data: the two terms of \eqref{eq:fr-affine} have norms $\|\bm M\|_2=3.9$ and $\|\bm K\|_2=6.7\times10^{11}$ (the stiffest element mode), so the perturbation term $|\theta_2(\omega)-\theta_2(\omega_j)|\,\|\bm K\|_2$ exceeds $\sigma_{\min}$ for any reference not at the frequency itself, and extra references on the same grid add nothing. The LP construction uses the same reference data far better and reaches \cfrSweepLpFinite{} of \cfrNGrid{} frequencies with \cfrSweepLpMax{} references, at a median overestimate of \cfrSweepLpRatio{}. Even then the certified bound is a statement that the error is below a number that is three orders of magnitude above it. A fifty-fold safety margin would be a useful bound; a thousand-fold margin tells the user little more than that the model has not diverged.

\begin{table}[ht]
\centering\small
\begin{tabular}{lrrr}
\toprule
bound & references & finite at (of 200) & median bound / true error \\
\midrule
Lipschitz & 15 & 1 & 3654 \\
Lipschitz & 40 & 1 & 3654 \\
Lipschitz & 100 & 1 & 3654 \\
Lipschitz & 200 & 1 & 3654 \\
LP & 15 & 8 & 4606 \\
LP & 40 & 23 & 3655 \\
LP & 80 & 50 & 1178 \\
\bottomrule
\end{tabular}

\caption{Frequencies (of \cfrNGrid{}) at which a finite certified bound exists, and the median ratio of the bound to the true error where it does, against the number of reference frequencies.}\label{tab:fr-sweep}
\end{table}

\begin{figure}[ht]
\centering
\includegraphics[width=0.98\linewidth]{figures/ch10_bound.pdf}
\caption{Left: true $\sigma_{\min}(\bm A(\omega))$ and the two lower bounds from $15$ references (bounds clipped at the plotting floor where they are zero). Right: true error and the certified bound where it is finite.}\label{fig:fr-bound}
\end{figure}

Cost of a query at this size: the Lipschitz bound takes \cfrQueryScmUs{}~$\mu$s, the LP bound \cfrQueryLpUs{}~$\mu$s, and an exact dense SVD of the $\cfrNFree\times\cfrNFree$ matrix \cfrQuerySvdUs{}~$\mu$s. At \cfrNFree{} unknowns there is therefore no saving from using a bound instead of the exact $\sigma_{\min}$; the construction pays off only when the full SVD is expensive, which was not tested.

\subsection{Random vibration}

The input is a flat PSD ($S_{\mathrm{in}}=1$) from $0.1\omega_1$ to $20\omega_1$ ($8.4$ to $1676$~Hz), applied through the tip pattern; the output is the tip displacement. The RMS of the full-order model is \cfrRvSig{} on a \cfrRvNPts{}-point grid. That grid matters: with $1500$ and $3000$ points the full-order RMS differs from the $6000$-point value by \cfrRvGridA and \cfrRvGridB, because the trapezoidal rule must resolve resonances whose half-power width at the first one is about $0.75$~Hz (estimated from the Rayleigh damping ratio; grid spacing $0.28$~Hz at $6000$ points). The comparison of \Cref{tab:fr-psd} is between the full-order and reduced PSD on the same $6000$-point grid, so it measures the basis and not the quadrature. Greedy bases with $5$ and $8$ columns agree with the full order to $10^{-10}$ in RMS; the $3$-column greedy basis does not (error $1.8\times10^{-2}$ in RMS and a maximum PSD error of $75\%$ of the peak), and a modal basis with only $3$ columns gives an RMS error of $8.6\times10^{-7}$. The RMS is dominated by the first resonance, so a basis that represents it well gives an accurate RMS even where the PSD at the other resonances is wrong; the maximum PSD error is the stricter measure. This test checks the reduced transfer function against the full one; \code{psd_response} and \code{FESystem.solve_random_vibration} implement the same formula, which was not checked against an independent closed form.

\begin{table}[ht]
\centering\small
\begin{tabular}{rlrr}
\toprule
columns & basis & RMS error & max PSD error / peak \\
\midrule
3 & greedy & $1.8\times10^{-2}$ & $7.5\times10^{-1}$ \\
3 & modal & $8.6\times10^{-7}$ & $1.4\times10^{-5}$ \\
5 & greedy & $1.7\times10^{-10}$ & $1.6\times10^{-9}$ \\
5 & modal & $2.4\times10^{-7}$ & $3.2\times10^{-6}$ \\
8 & greedy & $2.2\times10^{-10}$ & $2.5\times10^{-9}$ \\
8 & modal & $6.1\times10^{-8}$ & $7.8\times10^{-7}$ \\
\bottomrule
\end{tabular}

\caption{Random-vibration response of reduced models against the full-order result on the same $6000$-point grid: relative RMS error and the largest PSD error relative to the peak PSD.}\label{tab:fr-psd}
\end{table}

\section{Limits and failure modes}\label{sec:fr-limits}

\begin{limitsbox}
\begin{itemize}
\item \textbf{One model, one load, one band.} A \cfrNFree{}-unknown cantilever with one tip load and light Rayleigh damping. Rankings among the bases may differ for other structures, for non-proportional damping and for several loads.
\item \textbf{Greedy is not uniformly better.} It found the first two resonances and gave a lower median error at \cfrGreedyWinsMed{} of \cfrRankCount{} sizes, but a higher maximum error at four of the seven; at $5$ columns the maximum is set by an unvisited resonance.
\item \textbf{Indicators rank, they do not measure.} The residual norm ranks well here; the hierarchical indicator overestimates the current error by up to a factor of about \cfrHierOverMax{}.
\item \textbf{The certified bound is nearly vacuous on this model.} It is rigorous when finite, but it is finite at a handful of frequencies, never near a resonance, and at least three orders of magnitude above the true error where it exists. The Lipschitz construction does not improve with more references because $\|\bm K\|_2$ is large; a better-conditioned $\bm K$ was not tried, and the cost advantage of the bound at large $n$ was not measured.
\item \textbf{Random vibration needs a resolved grid} (\Cref{sec:fr-verif}). The RMS check says little about resonances above the first, and the formula was not checked against an independent closed form.
\end{itemize}
Not tested: the time-domain signal generator, the hierarchical indicator with a richer comparison model, an explicit non-proportional damping matrix, and models above \cfrNFree{} unknowns.
\end{limitsbox}

\section{Summary}\label{sec:fr-summary}

We started with five full-order solves and the question of where to place them. In the frequency domain the projection machinery of Chapter~\ref{ch:pod} applies with $\omega$ as the parameter, \eqref{eq:fr-affine}, and a snapshot basis is far more accurate than a modal one of the same size. The greedy loop finds the resonances unaided, but it is not uniformly better than an even grid: its median error is lower and its worst error is set by whichever resonance it has not yet visited. The residual norm is a good signal for where to add the next snapshot.

How far to trust the result is less settled. The certified bound \eqref{eq:fr-bound} is the one rigorous statement, and on this model it is finite at few frequencies and loose where it exists, so a test near the resonances, on frequencies not used for training, remains the working check. For a random load the reduced transfer function gives the RMS to $10^{-10}$ once the basis represents the first resonance and the frequency grid resolves the peaks. The next chapter leaves the frequency domain and reduces a structure by cutting it into components.

\section*{Exercises}\label{sec:fr-exercises}
\begin{enumerate}
\item \emph{(Implementation)} Repeat \Cref{tab:fr-greedy} with the third resonance inside the training band's seed set (for example with \code{n_seed=4}). Does the maximum error of the greedy basis drop below that of the uniform one?
\item \emph{(Implementation)} Run the certified bound on a model with $4$ elements instead of $20$ (so that $\|\bm K\|_2$ is much smaller). How many references does the Lipschitz construction need now, and what is the median overestimate?
\item \emph{(Implementation)} Replace the Rayleigh damping by an explicit non-proportional $\bm C$ (a single dashpot at the tip) and rebuild \cref{lst:fr-run} with the three-term decomposition. Is the greedy basis still accurate at $5$ columns?
\item \emph{(Implementation)} Double the number of frequencies in the PSD grid. Does the RMS agreement between the reduced and full models change, and does the full-order RMS?
\item \emph{(Implementation)} Time \code{residual_norm} against \code{hierarchical_error_indicator} for a basis of $10$ and of $40$ columns on a model of $2000$ unknowns. Which is cheaper, and by how much? (Test file: \texttt{rom\_engine/tests/test\_frequency.py}.)
\item \emph{(Derivation)} Let $\bm A(\omega)\bm x=\bm f$ and let $\bm x_r=\bm V\bm y$ be a reduced approximation with residual $\bm r=\bm f-\bm A\bm x_r$. Show that $\|\bm x-\bm x_r\|\le\|\bm r\|/\sigma_{\min}(\bm A(\omega))$, and say why the bound is large near a resonance.
\item \emph{(Interpretation)} The certified bound is nearly vacuous for the model of this chapter. State what a user may legitimately conclude from it, and what from the cheap indicator.
\end{enumerate}

\chapterrefs
